\section{Methodology}

We frame the music recommendation task as a zero-shot natural language generation problem. Given a textual representation of a user's listening history, the LLM is tasked with generating a list of 10 recommended tracks.

\subsection{Datasets}

\para{\musicrs}
We use the \musicrs benchmark \cite{surana2025musicrs}, which contains over 10,000 conversational music queries from Reddit. We extract the submission text as the "User Profile" and the comment entities (recommendations made by other human users) as the human-curated benchmark.

\para{\spotifydata}
To simulate real-world usage, we use a sample of raw Spotify listening history. Each session is represented as a sequence of Spotify Track IDs. We map these IDs to artist and track names using a metadata database to create a textual prompt for the LLM.

\subsection{Experimental Setup}

\para{LLM Configuration}
We evaluate two primary models: \claudesonnet and \gptfourmini. For the qualitative evaluation, we use \gptfour as an independent judge. All models are accessed via the OpenRouter API. We use a high temperature ($T=0.7$) to encourage diversity in the generated recommendations.

\para{Prompt Engineering}
The LLM is provided with a list of the user's recently played tracks (Track Name - Artist) and a brief instruction: "Based on this listening history, suggest 10 new tracks the user might enjoy. Provide a brief one-sentence explanation for each." This setup allows us to test both the recommendation quality and the model's reasoning capabilities.

\subsection{Evaluation Framework}

We employ two distinct evaluation strategies:

\para{Quantitative Hit Rate}
We measure the overlap between LLM-generated recommendations and the human-curated benchmark in the \musicrs dataset. A "hit" is defined as the LLM suggesting a track or artist that was also suggested by human commenters in the original Reddit thread.

\para{\judge}
To capture the nuances of music discovery, we conduct a side-by-side comparison. For a given user profile, we present the judge with two sets of recommendations (Set A and Set B, anonymized and randomized). One set is generated by the LLM, and the other is either a human-curated set or a baseline recommendation. The judge scores each set from 1--10 based on:
\begin{itemize}[leftmargin=*,itemsep=0pt,topsep=0pt]
    \item {\bf Relevance}: How well the tracks match the user's stated or implied preferences.
    \item {\bf Diversity}: The variety of genres and artists represented.
    \item {\bf Serendipity}: The degree to which the suggestions are non-obvious yet appropriate.
\end{itemize}

\subsection{Baselines}

We compare the LLMs against:
\begin{itemize}[leftmargin=*,itemsep=0pt,topsep=0pt]
    \item {\bf Human Curators}: The "gold standard" human recommendations from the \musicrs dataset.
    \item {\bf \bpr \cite{rendle2012bpr}}: A standard Bayesian Personalized Ranking model trained on a large-scale listening history dataset.
    \item {\bf \popularity}: A baseline that recommends the most popular tracks globally, used to measure the degree of personalization.
\end{itemize}
